\documentclass[11pt]{article}
\usepackage{amsmath,amssymb}
\usepackage[margin=1in]{geometry}

\title{Gravity Field Dynamics: Covariant Completion\\
  of the Constitutive Law\\[6pt]
  {\normalsize Supplementary Material for Letter~3}}
\author{Stephen Nelson\\Field Dynamics Research\\
  \texttt{stephen@fielddynamics.org}}
\date{}

\begin{document}
\maketitle

\begin{abstract}
This note provides the technical derivations referenced
in the accompanying letter.  Section~1
derives the unique kinetic functional $F(y)$ from the
constitutive law $\mu(x) = x/(x{+}1)$ by integration and
polynomial long division, confirming its three-term
structure.  Section~2 identifies every numerical prefactor
in the Gravity Field Dynamics (GFD) action with the
coupling polynomial $f(K,s) = 1 + sK + K^2$.  Section~3
derives the conformal electromagnetic coupling, shows that
$\gamma = 1$ follows from the connection structure (not
from the coupling itself), and identifies the coupling's
physical role as cosmological.  Section~4 varies the action and derives the
modified field equations, the spherically symmetric static
limit, and the PPN and gravitational-wave constraints.
\end{abstract}

% ============================================================
\section{From constitutive law to kinetic functional}
% ============================================================

The constitutive law $\mu(x) = x/(x{+}1)$, derived from
exchange symmetry between the two tetrad
roots~\cite{Nelson2026b}, determines the scalar-sector
kinetic functional through the AQUAL
identity~\cite{BekensteinMilgrom1984}:
\begin{equation}
  \mu(\sqrt{y}) = \frac{dF}{dy}\,,
  \label{eq:aqual_identity}
\end{equation}
where $y = |\nabla\Phi|^2/a_0^2$ is the squared
dimensionless acceleration.  Integration yields $F$
uniquely.

\subsection*{The integration}

Starting from the definition:
\begin{equation}
  F(y) = \int_0^y \mu(\sqrt{t})\,dt.
  \label{eq:F_integral}
\end{equation}
Substitute $u = \sqrt{t}$, so $t = u^2$ and $dt = 2u\,du$.
The limits become $u = 0$ to $u = \sqrt{y}$:
\begin{equation}
  F(y) = 2\int_0^{\sqrt{y}} \frac{u^2}{u + 1}\,du.
  \label{eq:F_substituted}
\end{equation}

Perform polynomial long division on the integrand:
\begin{equation}
  \frac{u^2}{u + 1} = u - 1 + \frac{1}{u + 1}.
  \label{eq:long_division}
\end{equation}
Verification: $(u - 1)(u + 1) + 1 = u^2 - 1 + 1 = u^2$.
\checkmark

Integrate term by term:
\begin{align}
  2\int_0^{\sqrt{y}} u\,du
    &= 2\left[\tfrac{u^2}{2}\right]_0^{\sqrt{y}}
    = y,
    \label{eq:term_K2} \\[4pt]
  2\int_0^{\sqrt{y}} (-1)\,du
    &= 2\bigl[-u\bigr]_0^{\sqrt{y}}
    = -2\sqrt{y},
    \label{eq:term_sK} \\[4pt]
  2\int_0^{\sqrt{y}} \frac{1}{u+1}\,du
    &= 2\bigl[\ln(u+1)\bigr]_0^{\sqrt{y}}
    = 2\ln(1 + \sqrt{y}).
    \label{eq:term_K0}
\end{align}
The result is the unique kinetic functional:
\begin{equation}
  \boxed{F(y) = y - 2\sqrt{y} + 2\ln(1 + \sqrt{y}).}
  \label{eq:F}
\end{equation}

\subsection*{Structural identification}

The three quotient terms in Eq.~(\ref{eq:long_division})
map to the three structural levels of the coupling
polynomial $f(K,s) = 1 + sK + K^2$
derived in Letter~2~\cite{Nelson2026b}:

\begin{center}
\renewcommand{\arraystretch}{1.3}
\begin{tabular}{lcll}
\hline
Quotient term & Integral & $F(y)$ term & Polynomial level \\
\hline
$u$ & $\int u\,du$ & $y$ & $K^2$ (face coupling) \\
$-1$ & $\int(-1)\,du$ & $-2\sqrt{y}$
  & $sK$ (propagation) \\
$1/(u{+}1)$ & $\int 1/(u{+}1)\,du$
  & $2\ln(1{+}\sqrt{y})$ & $K^0$ (field origin) \\
\hline
\end{tabular}
\end{center}

This identification is not pattern-matching.  The long
division decomposes the integrand into three terms whose
degrees match the three ranks of the tetrad equation.
The remainder $1/(u+1)$ equals $1 - \mu(u)$: the
logarithmic term integrates twice the TetB field fraction.
The polynomial enters the theory through the tetrad
structure (Letter~2) and re-emerges through the
variational structure (this integration).  Two independent
paths confirm the same three levels.

\subsection*{Limit checks}

At high acceleration ($y \gg 1$): $F(y)/y \to 1$.  The
scalar sector reduces to a standard kinetic term.  The
action becomes general relativity.

At low acceleration ($y \ll 1$): expanding
$\ln(1+\sqrt{y}) \approx \sqrt{y} - y/2 + y^{3/2}/3
- \cdots$ and collecting terms gives
$F(y) \to \tfrac{2}{3}\,y^{3/2}$.  This nonlinear kinetic
term produces the deep-field behaviour
$g_t \to \sqrt{g_s\,a_0}$, recovering the baryonic
Tully-Fisher relation $M \propto v^4$.

\subsection*{Derivative round-trip}

Differentiating Eq.~(\ref{eq:F}):
\begin{equation}
  \frac{dF}{dy} = 1 - \frac{1}{\sqrt{y}}
    + \frac{1}{\sqrt{y}(1+\sqrt{y})}
  = 1 - \frac{1}{1+\sqrt{y}}
  = \frac{\sqrt{y}}{1+\sqrt{y}}
  = \mu(\sqrt{y}).
\end{equation}
The round-trip $\mu \to F \to dF/dy \to \mu$ closes.
\checkmark

% ============================================================
\section{The numerical prefactors}
% ============================================================

The Gravity Field Dynamics (GFD) covariant action is
written in terms of the standard Newtonian gravitational
potential $\Phi$, which appears in both general relativity
and GFD.  The potential $\Phi$ is not a new field; it is the
same scalar potential whose gradient gives the gravitational
acceleration.  The torsion sector's contribution enters
through the constitutive law $\mu(x)$, which modifies the
kinetic structure from the standard $\nabla^2\Phi$ to
$\nabla\!\cdot[\mu\,\nabla\Phi]$.  The scalar sector of the
action is therefore not an additional field added to general
relativity; it is the energy functional of the
torsion-modified gravitational response, with the torsion
fraction $1 - \mu$ carried implicitly inside the kinetic
function $F(y)$ derived in Section~1.

The action is:
\begin{align}
  S_{\text{ST}} &= \int\!d^4x\,\sqrt{-g}\biggl[
    \frac{R}{16\pi G}
    - \frac{a_0^2}{8\pi G}\,F(y)
    \notag\\
    &\qquad\qquad
    -\; \tfrac{1}{4}\,e^{-2\Phi/c^2}
    F_{\mu\nu}F^{\mu\nu}
    + \mathcal{L}_{\text{matter}}
  \biggr].
  \label{eq:action}
\end{align}
Every numerical prefactor traces to the coupling
polynomial evaluated at $K = 4$ and $D = 3$.

\subsection*{Term 1: $R/(16\pi G)$}

The integer 16 is $K^2 = 4^2$: the number of
face-coupling channels in the $K \times K$ connection
matrix.  The solid angle per face is $4\pi/K = \pi$.
The gravitational constant $G$ is the measured CODATA
value; its internal structure in terms of the coupling
polynomial is derived
elsewhere~\cite{Nelson2026a,Nelson2026b}.  The prefactor
$16\pi G = K^2 \times \pi \times G$ compresses the
face-coupling count and the solid angle into the
curvature term.

\subsection*{Term 2: $a_0^2/(8\pi G)\,F(y)$}

The integer 8 is $2^D = 2^3$: the vertex-space
cardinality of the stellated octahedron (the number of
vertices of the cube that circumscribes the dual
tetrahedral pair).  The ratio of the two gravitational
prefactors is $16/8 = 2$, distinguishing single-chamber
($K^2$) from dual-chamber ($2^D$) quantities.

\subsection*{Term 3: $\tfrac{1}{4}\,e^{-2\Phi/c^2}
  F_{\mu\nu}F^{\mu\nu}$}

The integer 4 is $K$: the number of coupling faces per
tetrahedral chamber.  The three integers form a geometric
sequence:
\begin{equation}
  16,\; 8,\; 4 \;=\; K^2,\; 2^D,\; K,
  \qquad \text{common ratio } 2.
  \label{eq:geometric}
\end{equation}
Each successive term halves the structural count, moving
from the bilinear face coupling ($K^2$), through the
vertex-space dual structure ($2^D$), to the single-face
electromagnetic coupling ($K$).

\begin{center}
\renewcommand{\arraystretch}{1.3}
\begin{tabular}{llcl}
\hline
Action term & Prefactor & Integer & Polynomial origin \\
\hline
Curvature & $16\pi G$ & $K^2 = 16$
  & Face-coupling channels \\
Scalar sector & $8\pi G$ & $2^D = 8$
  & Vertex-space cardinality \\
EM coupling & $1/4$ & $K = 4$
  & Coupling faces per chamber \\
\hline
\end{tabular}
\end{center}

No prefactor is fitted.  Each is a structural count
from the stellated octahedron at $K = 4$, $D = 3$.

% ============================================================
\section{Gravitational lensing}
% ============================================================

\subsection*{Why lensing was thought to be broken}

Every scalar-tensor formulation of MOND since Bekenstein
\& Milgrom~\cite{BekensteinMilgrom1984} has faced a
lensing deficit: the theory predicts correct rotation
curves but too little photon deflection.
Bekenstein~\cite{Bekenstein2004} traced this to the PPN
parameter $\gamma$, which governs the ratio of spatial to
temporal metric potentials.  General relativity has
$\gamma = 1$, giving the observed deflection
\begin{equation}
  \theta = (1 + \gamma)\,\frac{2GM}{rc^2}
  = \frac{4GM}{rc^2}\,.
  \label{eq:deflection}
\end{equation}
In the scalar-tensor embedding of AQUAL, the modification
of gravity is carried by a scalar field $\phi$ that is
\textit{separate} from the metric.  This separate field
has its own stress-energy tensor, which contains
anisotropic stress (the
$\partial_\mu\phi\,\partial_\nu\phi$ terms).  The
anisotropy sources a difference between the temporal
potential $\Phi$ and the spatial potential $\Psi$, giving
$\gamma = \Psi/\Phi = \mu(x)$ in the quasi-static
limit.  At the galactic rim ($x \ll 1$), $\gamma \to 0$
and lensing drops to half the general-relativistic value.
This deficit motivated the construction of TeVeS, which
adds a timelike vector field to restore $\gamma = 1$ by
disformal coupling~\cite{Bekenstein2004}.

\subsection*{Why the deficit is a decomposition artefact}

The lensing deficit arises from a specific
\textit{embedding choice}: treating the constitutive
law's contribution as a separate scalar field with its
own stress-energy, bolted onto the 40-component
Christoffel connection.  In that embedding, the modified
Poisson equation is artificially decomposed as
\begin{equation}
  \underbrace{\nabla^2\Phi_N = 4\pi G\rho}_{\text{Newtonian (40 components)}}
  \;+\;
  \underbrace{\nabla\!\cdot\bigl[(\mu{-}1)\nabla\Phi\bigr]
  = 0}_{\text{scalar supplement}},
  \label{eq:decomposition}
\end{equation}
and the second equation is assigned its own
stress-energy tensor.  The anisotropy in that tensor is
what breaks $\gamma = 1$.

GFD does not perform this decomposition.  The
constitutive law $\mu(x)$ is not a separate field; it is
the partition of the full 64-component affine connection,
living inside the divergence operator.  The single
field equation is
\begin{equation}
  \nabla\!\cdot\bigl[\mu(x)\,\nabla\Phi\bigr]
  = 4\pi G\rho,
  \label{eq:gfd_poisson}
\end{equation}
with one potential $\Phi$ that governs the metric:
\begin{equation}
  g_{00} = -(1 + 2\Phi/c^2), \qquad
  g_{ij} = (1 - 2\Phi/c^2)\,\delta_{ij}.
  \label{eq:metric}
\end{equation}
Because the same $\Phi$ appears in both components,
$\gamma = \Psi/\Phi = 1$ at every radius and every
acceleration.

The logic is the same as in general relativity.  GR
has one field equation ($\nabla^2\Phi = 4\pi G\rho$),
one potential, and $\gamma = 1$.  GFD has one field
equation ($\nabla\!\cdot[\mu\,\nabla\Phi] = 4\pi G\rho$),
one potential, and $\gamma = 1$.  The only difference is
\textit{what} $\Phi$ satisfies.  GR freezes $\mu$ at~1.
GFD lets $\mu$ vary.  Neither introduces a second
potential.

\subsection*{Three frameworks compared}

\begin{center}
\renewcommand{\arraystretch}{1.3}
\begin{tabular}{lccc}
\hline
  Framework & Components & Rotation & $\gamma$ \\
\hline
GR (Christoffel)
  & 40 of 64 & wrong & $1$ \\
AQUAL-in-GR ($+\,\phi$)
  & 40 $+$ scalar & correct & $\mu(x)$ \\
GFD (full connection)
  & 64 of 64 & correct & $1$ \\
\hline
\end{tabular}
\end{center}

\noindent GR obtains $\gamma = 1$ but the wrong dynamics.
The scalar-tensor embedding obtains the correct dynamics
but $\gamma = \mu(x)$.  GFD obtains both: the correct
dynamics \textit{and} $\gamma = 1$, because the
modification lives inside the connection rather than
beside it.

\subsection*{Observational constraints}

The Cassini spacecraft measured
$\gamma - 1 = (2.1 \pm 2.3) \times
10^{-5}$~\cite{Bertotti2003}.  In GFD, $\gamma = 1$ at
all scales, satisfying this bound identically.
Gravitational waves propagate at $c$ because the tensor
sector is unmodified~\cite{Abbott2017}.

\subsection*{The electromagnetic coupling}

The electromagnetic sector of the
action~(\ref{eq:action}) carries a conformal coupling
to the scalar potential:
\begin{equation}
  \mathcal{L}_{\text{EM}}
    = -\tfrac{1}{4}\,e^{-2\Phi/c^2}\,
    F_{\mu\nu}F^{\mu\nu}.
  \label{eq:EM}
\end{equation}
Because $\gamma = 1$ is already guaranteed by the
connection-level formulation
(Eq.~\ref{eq:gfd_poisson}--\ref{eq:metric}), the
role of this coupling is not to repair a lensing deficit.
A conformal rescaling of the metric leaves null geodesics
invariant, so the coupling does not alter local photon
trajectories, and in GFD, no alteration is needed.
The term's physical role is cosmological: over
cosmological distances, the evolution of the background
potential $\Phi(z)$ modulates the electromagnetic action,
imprinting the gravitational scalar sector onto
photometric observables.  This produces a testable
redshift dependence in electromagnetic scaling relations
that is absent in the standard $\Lambda$CDM framework.

% ============================================================
\section{Variational derivation}
% ============================================================

\subsection*{The scalar field equation}

Varying the action~(\ref{eq:action}) with respect to the
scalar field $\Phi$ in the static limit, the scalar-sector
contribution is:
\begin{equation}
  \frac{\delta}{\delta\Phi}\left[
    -\frac{a_0^2}{8\pi G}\,F(y)
  \right]
  = -\frac{a_0^2}{8\pi G}\,\frac{dF}{dy}\,
    \frac{\partial y}{\partial\Phi}\,.
\end{equation}
Since $y = |\nabla\Phi|^2/a_0^2$ and
$dF/dy = \mu(\sqrt{y}) = \mu(|\nabla\Phi|/a_0)$,
the functional derivative yields a divergence:
\begin{equation}
  -\frac{1}{4\pi G}\,
  \nabla\!\cdot\!\bigl[
    \mu(|\nabla\Phi|/a_0)\,\nabla\Phi
  \bigr].
\end{equation}
Setting this equal to the matter source $-\rho$ (from
varying $\mathcal{L}_{\text{matter}}$) gives the
modified Poisson equation:
\begin{equation}
  \boxed{
  \nabla\!\cdot\!\bigl[
    \mu(|\nabla\Phi|/a_0)\,\nabla\Phi
  \bigr] = 4\pi G\rho.}
  \label{eq:modified_poisson}
\end{equation}
This is the Bekenstein--Milgrom
equation~\cite{BekensteinMilgrom1984} with the
constitutive law $\mu(x) = x/(x{+}1)$ fixed by the
frame geometry rather than fitted to data.
The electromagnetic coupling term contributes an
additional source proportional to the EM energy density,
which is negligible for non-relativistic matter
distributions.

\subsection*{Expansion of the modified Poisson equation}

Expanding the divergence:
\begin{equation}
  \mu(x)\,\nabla^2\Phi
  + (\nabla\mu)\cdot(\nabla\Phi) = 4\pi G\rho,
  \label{eq:expanded}
\end{equation}
where $x = |\nabla\Phi|/a_0$.  At high acceleration
($\mu \to 1$, $\nabla\mu \to 0$), this reduces to the
standard Poisson equation $\nabla^2\Phi = 4\pi G\rho$.
At low acceleration ($\mu \to 0$), the gradient term
$(\nabla\mu)\cdot(\nabla\Phi)$ dominates, producing the
nonlinear amplification that flattens rotation curves.

\subsection*{The spherically symmetric static limit}

In spherical symmetry, $|\nabla\Phi| = g_t$ (the total
gravitational acceleration governing circular orbits) and
the Poisson source gives $g_s = GM({<}r)/r^2$.
Eq.~(\ref{eq:modified_poisson}) reduces to the algebraic
relation:
\begin{equation}
  \mu(g_t/a_0)\,g_t = g_s.
  \label{eq:algebraic}
\end{equation}
Substituting $\mu(x) = x/(x{+}1)$:
\begin{equation}
  \frac{g_t}{g_t + a_0}\,g_t = g_s
  \qquad\Longrightarrow\qquad
  \boxed{g_s = \frac{g_t^2}{a_0 + g_t}\,.}
  \label{eq:sst}
\end{equation}
This is the scalar-tensor transition in the spherically
symmetric static limit.  At high acceleration
($g_t \gg a_0$): $g_s \to g_t$, Newtonian gravity.  At
low acceleration ($g_t \ll a_0$):
$g_t \to \sqrt{g_s\,a_0}$, the deep-field regime.

\subsection*{PPN parameters}

In the high-acceleration limit $\mu \to 1$, the kinetic
functional $F(y)/y \to 1$ (Section~1).  The scalar sector
becomes a standard minimally coupled kinetic term and
decouples from the metric dynamics.  The action reduces
to Einstein--Hilbert plus a decoupled scalar plus
standard matter.  All post-Newtonian parameters inherit
their general-relativistic values:
$\beta = \gamma = 1$.  As shown in Section~3, $\gamma = 1$
holds at all accelerations because the connection-level
formulation produces a single metric potential.  The Cassini
bound~\cite{Bertotti2003} is satisfied.

\subsection*{Gravitational wave speed}

Gravitational waves are tensor perturbations $h_{\mu\nu}$
of the metric.  The scalar sector of the
action~(\ref{eq:action}) modifies only the scalar
potential $\Phi$; it does not couple to the tensor
degrees of freedom.  The tensor sector is the unmodified
Einstein--Hilbert term $R/(16\pi G)$, whose linearisation
gives the standard wave equation
$\Box\,h_{\mu\nu} = 0$.  Gravitational waves therefore
propagate at $c$ on the background metric.  The
multi-messenger observation of GW170817 constrains
$|c_{\text{gw}}/c - 1| < 10^{-15}$~\cite{Abbott2017}.
The GFD action satisfies this bound identically.

\subsection*{Summary of field equations}

The complete downstream structure of the GFD action in
the non-relativistic, spherically symmetric limit:
\begin{align}
  &\text{Modified Poisson:} &
    &\nabla\!\cdot\!\bigl[
      \mu(|\nabla\Phi|/a_0)\,\nabla\Phi
    \bigr] = 4\pi G\rho,
    \label{eq:summary_poisson} \\[4pt]
  &\text{SST (spherical):} &
    &g_s = g_t^2/(a_0 + g_t),
    \label{eq:summary_sst} \\[4pt]
  &\text{PPN:} &
    &\beta = \gamma = 1,
    \label{eq:summary_ppn} \\[4pt]
  &\text{GW speed:} &
    &c_{\text{gw}} = c.
    \label{eq:summary_gw}
\end{align}
No free functions.  No free parameters beyond $G$ and
$a_0$.  The constitutive law $\mu(x) = x/(x{+}1)$
determines the action; the action determines the field
equations; the field equations recover general relativity
at high acceleration and the empirical radial
acceleration relation at low acceleration.

% ============================================================

\begin{thebibliography}{99}

\bibitem{Nelson2026a}
S.~Nelson,
\textit{The Coordinate-Space Geometry of the Tetrad
Frame},
Class.\ Quantum Grav.\ (2026), submitted (CQG-115428).

\bibitem{Nelson2026b}
S.~Nelson,
\textit{The Other Half of Gravity: Torsion and the
Two-Chamber Structure of the Tetrad Frame},
Class.\ Quantum Grav.\ (2026), submitted (CQG-115447).

\bibitem{BekensteinMilgrom1984}
J.~Bekenstein and M.~Milgrom,
\textit{Does the missing mass problem signal the breakdown
of Newtonian gravity?},
Astrophys.\ J.\ \textbf{286} (1984) 7.

\bibitem{Bekenstein2004}
J.~D.~Bekenstein,
\textit{Relativistic gravitation theory for the modified
Newtonian dynamics paradigm},
Phys.\ Rev.\ D \textbf{70} (2004) 083509.

\bibitem{Bertotti2003}
B.~Bertotti, L.~Iess and P.~Tortora,
\textit{A test of general relativity using radio links
with the Cassini spacecraft},
Nature \textbf{425} (2003) 374.

\bibitem{Abbott2017}
B.~P.~Abbott \textit{et al.}\ (LIGO/Virgo and
Fermi GBM Collaborations),
\textit{Gravitational Waves and Gamma-Rays from a Binary
Neutron Star Merger: GW170817 and GRB~170817A},
Astrophys.\ J.\ Lett.\ \textbf{848} (2017) L13.

\end{thebibliography}

\end{document}
